\section{Putting it all together}\label{sec:wrap}

\subsection{The story of the numbers}

Follow one Fashion-MNIST MLP picture through the whole project:

\begin{center}
\small
\begin{tabular}{lrrl}
\toprule
Version & million cycles & vs.\ start & what changed\\
\midrule
Exact, one call per multiply (paper style) & 16.04 & 1.00$\times$ & starting point\\
RoBA-Dn1, one call per multiply & 8.75 & 1.83$\times$ & LOD + approximate multiply\\
Exact, fused loop & 13.49 & 1.19$\times$ & better loop, exact multiply\\
RoBA-Dn1, fused loop & 7.05 & 2.28$\times$ & better loop + RoBA\\
Table of squares & 3.87 & 4.14$\times$ & no multiply at all, exact\\
Table of squares, raw pixels + skip & 2.70 & 5.94$\times$ & also skip the black background\\
\bottomrule
\end{tabular}
\end{center}

And the accuracy of all of them on the full test set stays between 89.5\% and 89.8\%.

\subsection{What we learned}

\begin{enumerate}
  \item \textbf{Approximate multiplication does not hurt these networks.} Every one of the seven
        approximate multipliers stays within 0.75 points of exact at int8, most within 0.2 points, even
        RoBA-Dn1, which is 9\% too small on average. Fine-tuning recovers most of the rest.
  \item \textbf{The LOD instruction is essential.} With it, Dn1 costs 28 cycles per multiply at any size;
        without it, 91 or more, which is slower than exact.
  \item \textbf{Faster multiplies expose the rest of the loop.} Per multiply Dn1 is 3$\times$ faster, per
        picture only 1.8$\times$, because the wrapper and the loop then dominate (Amdahl's law). Fusing the
        multiply into the loop raises Dn1's advantage to 1.9--2.4$\times$.
  \item \textbf{More accurate is not always better on this CPU.} The two-step RoBA methods and DRUM6 are
        slower than exact for 8-bit numbers and use more energy. Only the one-step methods (Dn1, UpDn,
        Mitchell) and Dn2 pay off.
  \item \textbf{For int8, an exact table of squares wins.} It needs 1~KB and no new instruction, gives exact
        products, and is 3.0--3.6$\times$ faster than exact (1.3--1.8$\times$ faster than Dn1).
  \item \textbf{Zeros matter.} Feeding raw pixels (so the background is 0) and skipping zeros brings the MNIST
        MLP to 10.5$\times$.
  \item \textbf{Energy follows cycles.} LOD adds 4.5\% power, so Dn1 at about 2$\times$ the speed uses roughly
        half the energy, and the table about 0.3.
  \item \textbf{It holds on real tasks, and precision picks the winner} (Phase 3). On CIFAR-10, keyword spotting
        and machine-sound anomaly detection, with three seeds: at int8 the table wins (2.0--2.6$\times$, exact);
        at int16 the table does not fit and RoBA wins (Dn1 2.9--3.3$\times$ at a third of the energy, Dn2
        2.4--2.5$\times$ within 0.2 points of exact); at int4 nothing gains much.
  \item \textbf{It holds on a real chip, and memory speed matters} (Phase 4). On the Zybo FPGA the Ibex gives the
        same outputs as the Python model (960/960) and, from on-chip RAM, exactly the simulated cycles (240/240).
        The LOD costs no measurable area. With data in DDR (about 15 cycles per access) the table drops to about
        exact speed, while RoBA-Dn1 keeps 1.3--1.4$\times$ at int8 and 2.0--2.2$\times$ at int16.
\end{enumerate}

\subsection{Limits: what these results do not show}

\begin{itemize}
  \item \textbf{One training seed} per dataset in Phases 1--2 (Phase 3 uses three): there, accuracy differences
        below about 0.1 points are noise.
  \item \textbf{100 test pictures} on Ibex per configuration (the full 10{,}000 are covered by the bit-exact
        Python model). The cycle counts vary very little from picture to picture, so 100 is enough for them.
  \item \textbf{Energy is estimated}, not measured: cycles times relative power from the paper. The FPGA's
        Vivado power report is itself an estimate (default switching rates, 1~mW resolution); no current was
        measured on the board.
  \item \textbf{One-cycle memory} in the simulator. The FPGA (Phase 4) measured the other extreme, DDR at about
        15 cycles per access; caches or flash wait states would lie in between.
  \item \textbf{Phase 3 cycles use seed 0 and 20 inputs} per configuration; accuracy uses all seeds and the
        full test sets through the bit-exact Python model.
  \item \textbf{Keyword spotting at int4 is unsolved} (about 61\% for every method); keeping the first layer at
        8 bits is the likely fix and has not been tried.
  \item \textbf{We did not measure a core with a hardware multiplier}; it would need fewer cycles, at more area
        and power.
  \item Our \textbf{Mitchell} is the classic one with the carry; the paper's Mitchell gives the same products as
        RoBA-Dn1.
\end{itemize}

\subsection{What comes next}
\begin{enumerate}
  \item The two remaining datasets in the Phase 3 pipeline: Visual Wake Words (MLPerf Tiny's person detection)
        and MNIST / Fashion-MNIST with three seeds.
  \item Fix int4 keyword spotting (first layer at 8 bits).
  \item Choosing the multiplier per layer: cheap where the network tolerates errors, accurate where it is
        sensitive.
  \item Replace the equal-power assumption by power per method: Vivado power analysis with each program's
        simulated switching activity, or a current measurement on the board.
  \item A small data cache or tighter memory layout, so that more programs run from on-chip RAM.
\end{enumerate}

\subsection{Glossary}

\begin{description}[style=nextline,leftmargin=1.2cm,font=\normalfont\bfseries]
  \item[Accumulator (\code{acc})] The running sum inside a neuron while its products are added up.
  \item[Activation] A value flowing between layers (the input of a neuron).
  \item[ALU] Arithmetic logic unit: the part of the CPU that adds, shifts, compares, and now also does LOD.
  \item[Amdahl's law] Speeding up one part of a job by $k$ only helps in proportion to that part's share of
        the time.
  \item[Approximate square] $\tilde S(v) = (2v - 2^k)\,2^k$ with $k = \mathrm{LOD}(v)$; a cheap but inexact
        square. Fails inside the polarization identity.
  \item[Approximate multiplier] A method that computes a product that is close to, but not always exactly,
        the true product, in exchange for speed or energy.
  \item[AUC] Area under the ROC curve: the chance that a random broken example scores higher than a random
        normal one. 50\% is coin tossing, 100\% is perfect.
  \item[Autoencoder] A network that squeezes its input through a narrow middle and rebuilds it; a bad rebuild
        means an unusual input.
  \item[Bare metal] A program that runs directly on the CPU, with no operating system.
  \item[Bitstream] The file that configures an FPGA into a particular circuit.
  \item[Bias] The extra number a neuron adds after its sum of products.
  \item[Block RAM] Memory blocks inside the FPGA (4~KB each, 60 on the XC7Z010); one-cycle access.
  \item[Bit, binary] A 0/1 digit; numbers written with place values 1, 2, 4, 8, \dots
  \item[Calibration] Choosing the scale that maps each layer's decimal values onto the integer range.
  \item[Checkpoint] A saved file with all the weights of a trained network.
  \item[Clock cycle] One tick of the CPU's clock; the unit we measure speed in.
  \item[Convolution] A layer that slides a small filter of weights over a picture.
  \item[Custom-0] An opcode RISC-V reserves for non-standard instructions; LOD lives there.
  \item[Cycle-accurate simulation] A simulation that reproduces the chip tick by tick, so its cycle counts
        are exact.
  \item[Depthwise-separable convolution] A $3\times3$ convolution per channel followed by a $1\times1$
        convolution across channels; about 8 times fewer multiplications. Used by DS-CNN.
  \item[DDR] The board's main memory chip: large (512~MB) but slow for the Ibex (about 15 cycles per access).
  \item[Decoder] The part of the CPU that works out what an instruction means.
  \item[DRUM6] An approximate multiplier that keeps the top 6 bits of each number and multiplies those
        exactly.
  \item[Epoch] One pass over all training pictures.
  \item[Fine-tuning] A few more training epochs, here with the approximate multiplier switched on.
  \item[Fixed point] Representing a decimal as an integer with an agreed scale (here $2^{24}$).
  \item[FPGA] A chip whose circuit can be reprogrammed; used to try out CPU designs for real.
  \item[FPGA board, Zybo] Digilent's board with a Zynq XC7Z010 (ARM processor + FPGA) used in Phase 4.
  \item[Fused MAC] Writing the multiplier directly inside the multiply-accumulate loop instead of calling it.
  \item[Ibex] A small open-source RISC-V CPU from lowRISC.
  \item[int4 / int8 / int16] Whole numbers stored in 4, 8 or 16 bits ($\pm7$, $\pm127$ or $\pm32767$ for weights).
  \item[LeNet-5] A small convolutional network: two convolution layers and three fully connected layers.
  \item[LOD] Leading-one detector: the instruction that returns the position of a number's highest 1 bit.
  \item[Logit / score] A raw output value of the last layer; the largest one is the answer.
  \item[Logit hash] A 32-bit fingerprint of the 10 scores, used to check them quickly.
  \item[Loss] A number that measures how wrong the network is; training makes it small.
  \item[M extension] The optional RISC-V instructions for multiplication and division; our core does not have
        them.
  \item[MFCC] Mel-frequency cepstral coefficients: a small ``picture of a sound'' (energy at different pitches
        over time).
  \item[MLPerf Tiny] The standard benchmark of tasks and reference networks for machine learning on tiny chips.
  \item[LUT] Look-up table: the FPGA's basic logic block; area is counted in LUTs and flip-flops.
  \item[MAC] Multiply-accumulate: one product added to a running sum.
  \item[Mitchell] A classic approximate multiplier based on logarithms.
  \item[MLP] Multi-layer perceptron: a network of fully connected layers.
  \item[MNIST, Fashion-MNIST] Datasets of $28\times28$ grey pictures: hand-written digits and clothes.
  \item[Normalisation] Shifting and scaling the pixels to be centred around 0 ($(x-\text{mean})/\text{std}$).
  \item[Polarization identity] $xw = ((x+w)^2 - x^2 - w^2)/2$: a product from three squares.
  \item[Priority encoder] A circuit that outputs the position of the highest 1 among its inputs.
  \item[QAT] Quantization-aware training: fine-tuning with the integer rounding (and here the approximate
        multiplier) switched on in the forward pass.
  \item[Quantization] Turning decimal numbers into integers with a scale.
  \item[Raw pixels] Feeding the picture without normalisation (the normalisation is folded into the first
        layer), so black is 0.
  \item[ReLU] The rule ``negative $\to$ 0, positive stays''.
  \item[ResNet] A network with shortcuts that add a block's input to its output.
  \item[Register] One of the CPU's 32 small storage slots for numbers.
  \item[RISC-V] An open instruction set; RV32I is its basic 32-bit integer version, C adds short
        instructions.
  \item[RoBA] Rounding-based approximate multiplication: round to powers of two, keep the cheap parts.
        Dn = round down, UpDn = round to the nearest, the ``2'' versions repeat it on the leftover.
  \item[Seed] The starting number of the random-number generator; a different seed gives a different
        training run.
  \item[Shift] Moving all bits left (multiply by 2) or right (divide by 2).
  \item[Shift-and-add] The exact way to multiply without a multiplier: add shifted copies.
  \item[Sign-magnitude] Storing a number as a sign bit plus its size.
  \item[Square table] A 509-entry table of $v^2$ used by the polarization identity.
  \item[Straight-through estimator] Training trick: use the approximate result going forward, pretend it
        was exact going backward.
  \item[Test accuracy] Share of unseen test pictures the network gets right.
  \item[Verilator] A tool that turns a hardware description into a fast C++ cycle-accurate simulator.
  \item[Vivado, Vitis] Xilinx tools: hardware to bitstream (with area/timing/power reports); ARM software.
  \item[WNS] Worst negative slack: the timing margin of the slowest path; positive means the clock is safe.
  \item[Zero skipping] Not doing the work for inputs that are 0.
\end{description}
\newpage
